% vim: ai sw=2

\documentclass[xcolor=dvipsnames,handout]{beamer}

\let\Oldemph\emph
\renewcommand{\emph}[1]{\textcolor{blue}{\Oldemph{#1}}}
\newcommand{\heading}[1]{\textbf{\textcolor{BurntOrange}{#1}}}
\newcommand{\header}[1]{\textbf{\texttt{\textcolor{WildStrawberry}{#1}}}}
\newcommand{\doc}[1]{\textbf{\texttt{\textcolor{OliveGreen}{#1}}}}

% Copyright 2004 by Till Tantau <tantau@users.sourceforge.net>.
%
% In principle, this file can be redistributed and/or modified under
% the terms of the GNU Public License, version 2.
%
% However, this file is supposed to be a template to be modified
% for your own needs. For this reason, if you use this file as a
% template and not specifically distribute it as part of a another
% package/program, I grant the extra permission to freely copy and
% modify this file as you see fit and even to delete this copyright
% notice. 


\mode<presentation>
{
%  \usetheme{Warsaw}
\usetheme{Madrid}

  \setbeamercovered{transparent}
}

\usepackage[czech]{babel}
\usepackage[utf8]{inputenc}
\usepackage{times}
\usepackage[IL2]{fontenc}

\usepackage{graphicx}
\usepackage[final]{listings}
\usepackage{multirow}

\lstset{tabsize=2,numbers=none,showstringspaces=false,breaklines=true,
	breakatwhitespace=true,escapeinside={(*@}{@*)},
	basicstyle=\scriptsize,keywordstyle=\bfseries\color{OliveGreen},
	commentstyle=\color{blue}}

\title[PB173/01] % (optional, use only with long paper titles)
{PB173 -- Ovladače jádra -- Linux}

\subtitle{VI.}

\author[J. Slabý]{Jiří~Slabý}

\institute[ITI] % (optional, but mostly needed)
{ITI, Fakulta Informatiky}

\date{02. 11. 2010}

\subject{}
% This is only inserted into the PDF information catalog. Can be left
% out. 



% If you have a file called "university-logo-filename.xxx", where xxx
% is a graphic format that can be processed by latex or pdflatex,
% resp., then you can add a logo as follows:

% \pgfdeclareimage[height=0.5cm]{university-logo}{university-logo-filename}
% \logo{\pgfuseimage{university-logo}}



% Delete this, if you do not want the table of contents to pop up at
% the beginning of each subsection:
% \AtBeginSubsection[]
% {
%   \begin{frame}<beamer>
%     \frametitle{Outline}
%     \tableofcontents[currentsection,currentsubsection]
%   \end{frame}
% }


% If you wish to uncover everything in a step-wise fashion, uncomment
% the following command: 

%\beamerdefaultoverlayspecification{<+->}

\begin{document}

\begin{frame}
  \titlepage
\end{frame}

%%%%%%%%%%%%%%%%%%%%%%%

\begin{frame}
  \frametitle{Datové typy}

  \heading{LDD3 kap. 11}

  \bigskip

  \begin{itemize}
  \item Základní typy
%  \item Přenositelnost
  \item Endianita (pořadí bytů)
  \item Datové struktury (seznam, FIFO)
  \end{itemize}

\end{frame}

\begin{frame}
  \frametitle{Základní typy}

  \begin{itemize}
  \item \texttt{char}, \texttt{short}, \texttt{int}, \texttt{long},
    \texttt{long long} (+ \texttt{unsigned})
  \item \texttt{u8}, \texttt{u16}, \texttt{u32}, \texttt{u64},
    \texttt{\_\_u8}, \texttt{\_\_u16}, \texttt{\_\_u32}, \texttt{\_\_u64}
  \item \texttt{s8}, \texttt{s16}, \texttt{s32}, \texttt{s64},
    \texttt{\_\_s8}, \texttt{\_\_s16}, \texttt{\_\_s32}, \texttt{\_\_s64}
  \end{itemize}

  \begin{table}
  \footnotesize
  \begin{tabular}{|l|c|c|c|c|c|c|} \hline
arch    & char & short & int & long & ptr & long long \\ \hline\hline
alpha   & 1    & 2     & 4   & 8    & 8   & 8 \\ \hline
arm     & 1    & 2     & 4   & 4    & 4   & 8 \\ \hline
ia64    & 1    & 2     & 4   & 8    & 8   & 8 \\ \hline
m68k    & 1    & 2     & 4   & 4    & 4   & 8 \\ \hline
mips    & 1    & 2     & 4   & 4    & 4   & 8 \\ \hline
ppc     & 1    & 2     & 4   & 4    & 4   & 8 \\ \hline
sparc   & 1    & 2     & 4   & 4    & 4   & 8 \\ \hline
sparc64 & 1    & 2     & 4   & 8    & 8   & 8 \\ \hline
x86\_32 & 1    & 2     & 4   & 4    & 4   & 8 \\ \hline
x86\_64 & 1    & 2     & 4   & 8    & 8   & 8 \\ \hline
  \end{tabular}
  \caption{\texttt{sizeof} základních typů}
  \end{table}

\end{frame}

\begin{frame}
  \frametitle{Základní typy}

  \begin{table}
  \footnotesize
  \begin{tabular}{|l|c|c|c|} \hline
arch        & virt & phys  & long \\ \hline\hline
x86\_32     & 4    & 4     & 4    \\ \hline
x86\_32+PAE & 4    & 4.5   & 4    \\ \hline
x86\_64     & 8    & 8     & 8    \\ \hline
  \end{tabular}
  \caption{\texttt{sizeof} dalších typů}
  \end{table}

  \heading{Pro práci s fyzickou adresou \texttt{long} nestačí}
  \begin{itemize}
  \item Ale \texttt{u64} může být zbytečné (čisté 32-bity)
  \item Proto speciální typy
  \begin{itemize}
  \item \texttt{phys\_addr\_t} (RAM)
  \item \texttt{dma\_addr\_t} (RAM)
  \item \texttt{resource\_size\_t} (PCI prostor apod.)
  \item Anebo \texttt{long}, ale pak v něm uchovávat \texttt{pfn}
    (\texttt{phys/PAGE\_SIZE})
  \end{itemize}
  \end{itemize}

  \pause

  Úkol: opravte typy v pb173/06 (první TODO)

\end{frame}

\begin{frame}
  \frametitle{Pořadí bytů}

  \begin{table}
  \caption{Číslo 0x12345678 v paměti}
  \begin{tabular}{l|llll}
Adresa:       &    0 &    1 &    2 &    3 \\ \hline
Big-endian    & 0x12 & 0x34 & 0x56 & 0x78 \\
Little-endian & 0x78 & 0x56 & 0x34 & 0x12 \\
Mixed-endian  & 0x34 & 0x12 & 0x78 & 0x56 \\
  \end{tabular}
  \end{table}

  \pause
  \medskip

  \heading{Pro typy o velikosti větší než 1 je nutné vědet pořadí}
  \begin{itemize}
  \item Specifikace
  \begin{itemize}
  \item x86: LE
  \item PowerPC: LE nebo BE (dle nastavení CPU)
  \item PCI: LE
  \item Síťový provoz: BE
  \end{itemize}
  \end{itemize}

\end{frame}

\begin{frame}
  \frametitle{Funkce pro pořadí bytů}

  \begin{itemize}
  \item \header{asm/byteorder.h}
  \item \texttt{\_\_leXX}, \texttt{\_\_beXX}
  \item \texttt{cpu\_to\_leXX}, \texttt{cpu\_to\_beXX}
  \item \texttt{leXX\_to\_cpu}, \texttt{beXX\_to\_cpu}
  \item $\text{\texttt{XX}}\in\left\{\text{\texttt{16}},\text{\texttt{32}},\text{\texttt{64}}\right\}$
  \item \texttt{*p} varianty -- ukazatele
  \item \texttt{*s} varianty -- in-place
  \item \texttt{ntohs}, \texttt{ntohl}, \texttt{htons}, \texttt{htonl}
  \end{itemize}

  \pause
  \medskip

  Úkol: doplňte tělo \texttt{decode\_and\_print} v pb173/06 (druhé TODO)

\end{frame}

\begin{frame}
  \frametitle{Chyby jako ukazatele}

  \begin{itemize}
  \item \texttt{void *funkce()}
  \item Jak vrátit chytřejší chybu, než NULL?
  \item Díra na konci adresového prostoru
  \begin{itemize}
  \item Např. \texttt{0xfffffffffff00000-0xffffffffffffffff} (x86\_64)
  \item Lze použít k zakódování chyby
  \end{itemize}
  \item \texttt{void *ptr=ERR\_PTR(-Eerror)}
  \item \texttt{IS\_ERR(ptr)} $\Rightarrow$ \texttt{int err=PTR\_ERR(ptr)}
  \item \texttt{IS\_ERR(NULL)=false}
  \end{itemize}

\end{frame}

\begin{frame}[fragile]
  \frametitle{Seznam I.}

  \begin{itemize}
  \item \header{linux/list.h}
  \item \texttt{struct list\_head} (obsahuje \texttt{prev}, \texttt{next})
  \begin{itemize}
  \item Jak samotný seznam (počátek), tak jeho prvky
  \end{itemize}
  \item \texttt{LIST\_HEAD(my\_list)}, \texttt{INIT\_LIST\_HEAD(\&my\_list)}
  \item \texttt{list\_add(co, kam)}, \texttt{list\_add\_tail},
    \texttt{list\_del}
  \item \texttt{list\_move}, \texttt{list\_move\_tail}
  \end{itemize}

  \begin{center}
  \begin{tabular}{ll}
  \begin{lstlisting}[language=C]
struct my_struct {
  int a;
  struct list_head list;
};
LIST_HEAD(my_list);
  \end{lstlisting} &
  \begin{lstlisting}[language=C]
struct my_struct *s;
s = kmalloc(sizeof(*s), GFP_KERNEL);
INIT_LIST_HEAD(&s->list);
list_add(&s->list, &my_list);
/* list_add_tail(&s->list, &my_list); */
  \end{lstlisting}
  \end{tabular}
  \end{center}

  \pause
  \medskip

  Úkol: naalokujte 20 stránek a v seznamu si je pamatujte.

\end{frame}

\begin{frame}[fragile]
  \frametitle{Seznam II.}

  \begin{itemize}
  \item \header{linux/list.h}
  \item \texttt{list\_empty}
  \item \texttt{list\_first\_entry}, \\
    \texttt{list\_for\_each\_entry}, \\
    \texttt{list\_for\_each\_entry\_reverse}
  \item \texttt{*\_safe} varianty -- pokud měním aktuální člen seznamu
  \end{itemize}

  \begin{lstlisting}[language=C]
struct my_struct *s;
s = list_first_entry(&my_list, struct my_struct, list);
/* or */
list_for_each_entry(s, &my_list, list) {
  s->a;
}
  \end{lstlisting}

  \pause
  \medskip

  Úkol: předchozích 20 stránek uvolněte.

\end{frame}

\begin{frame}
  \frametitle{FIFO}

  \begin{itemize}
  \item \header{linux/kfifo.h}, \doc{samples/kfifo/*}
  \item \texttt{struct kfifo}
  \item \texttt{DECLARE\_KFIFO} $\Rightarrow$ \texttt{INIT\_KFIFO}
  \item Jinak: \texttt{kfifo\_alloc(\&fifo, size, GFP\_*)}, \texttt{kfifo\_free}
  \item Zápis: \texttt{kfifo\_in(\&fifo, buf, count)}
  \item Čtení: \texttt{kfifo\_out(\&fifo, buf, count)}
  \item Obsazeno: \texttt{kfifo\_len(\&fifo)}
  \item Volno: \texttt{kfifo\_avail(\&fifo)}
  \end{itemize}

\end{frame}

\begin{frame}
  \frametitle{Úkol}

  \heading{Práce s FIFO} (\header{linux/kfifo.h})
  \begin{enumerate}
  \item Globalní FIFO 8 intů (\texttt{DECLARE\_KFIFO})
  \item \texttt{INIT\_KFIFO}
  \item Zapisovat čísla (\texttt{kfifo\_in}) dokud je místo
    (\texttt{kfifo\_avail})
  \item Vypsat všechna čísla (\texttt{kfifo\_out} + \texttt{kfifo\_len})
  \end{enumerate}

\end{frame}

\end{document}

